\section{Conclusion}
We present an investigation of the use of an agent-based repair system tailored to Google’s internal software development environment. We collected a benchmark set from Google's internal issue tracking system.
To establish floor for agent-based APR performance on this benchmark, 
we developed Passerine, a minimal agentic repair system with access to Google development tools. We showed that Passerine’s minimal tooling and basic ReAct-style prompting can successfully generate (with 20 trajectory samples) 
plausible patches for 73\% and 25.6\% machine-reported
and human-reported bugs, respectively,
in our evaluation set. Manual inspection showed that 
43\% and 17.9\% of machine-reported bugs 
and human-reported bugs, respectively, had at least one valid patch semantically matching the ground truth. We also find that Passerine automatically adapts its behavior to take advantage of the varying amounts of information provided in machine- and human-reported bugs.
To place these results in context, we also studied how Google's environment can differ from those used in a popular open source repair benchmark, with much larger and spread out patches and with bug descriptions that contain fewer likely code symbols. 


\section{Acknowledgements}
We thank Sherry Shi, Sam Cheng, and Alex Polozov for their valuable feedback and assistance.
